% 第8章扩展实验建议新增内容
% 建议插入到“实验设计与复现设置”之后、“实验结果及分析”之前。

\subsection{扩展实验规模与受控合成基准}
\label{subsec:扩展实验规模与受控合成基准}

为进一步验证方法在不同模型规模和异常类型下的表现，本文在原有8个BPMN、67个日志输入的基础上构建了扩展实验基准。扩展基准由两部分组成：第一部分保留原有8个BPMN协作图，通过对其中能够被所有参与方接受的正常全局消息迹执行消息删除、相邻消息换序以及发送端/接收端可见性变异，新增64个日志输入；第二部分生成6个结构化、无循环的BPMN协作图，参与方数量由5增加到20，消息流数量由12增加到96，并为每个模型生成15个日志输入，共得到90个合成日志。扩展后，实验包含14个不同BPMN协作模型和221个日志输入。

所有新增数据均由脚本\path{scripts/generate_extended_benchmark.py}在固定随机种子20260809下生成。对于规模化合成模型，协调者子流程由顺序块、并行块和排他选择块组合而成，其余参与方通过消息流与协调者交互。每个模型分别注入5个消息缺失异常、5个顺序逆转异常和5个互斥分支同时执行异常。对于既有BPMN，生成器首先从原日志中筛选被全部参与方接受的正常迹，再对其执行候选变异。每个新增日志都包含生成种子、异常类型和注入位置注释；JSON和CSV清单进一步记录BPMN及日志的SHA-256、受影响消息、检出的参与方、MinAS表达式和通过严格复验的表7-5规则编号。

\begin{table}[htbp]
\centering
\small
\renewcommand{\arraystretch}{1.25}
\caption{扩展前后的实验数据集构成}
\label{tab:扩展实验数据集构成}
\begin{tabularx}{\textwidth}{
    >{\centering\arraybackslash\hsize=1.25\hsize}X
    >{\centering\arraybackslash\hsize=0.65\hsize}X
    >{\centering\arraybackslash\hsize=0.75\hsize}X
    >{\centering\arraybackslash\hsize=0.80\hsize}X
    >{\centering\arraybackslash\hsize=1.55\hsize}X
}
\toprule
数据组 & BPMN数量 & 日志输入数 & 模型规模范围（参与方、消息流） & 主要用途 \\
\midrule
原始人工实验 & 8 & 67 & $(2,4)$--$(9,25)$ & 多工业场景案例验证和局限性分析 \\
既有BPMN变异日志 & 复用8 & 64 & $(2,4)$--$(9,25)$ & 既有模型上的回归验证和异常变体覆盖 \\
规模化合成BPMN & 6 & 90 & $(5,12)$--$(20,96)$ & 受控正确性验证、异常类型覆盖和规模性能分析 \\
\midrule
合计 & 14个不同BPMN & 221 & $(2,4)$--$(20,96)$ & -- \\
\bottomrule
\end{tabularx}
\end{table}

新增154个日志的异常类型分布为：消息缺失51个、合成模型顺序逆转30个、互斥分支同时执行30个、既有模型相邻消息换序9个、仅发送端可见16个、仅接收端可见18个。独立验证脚本\path{scripts/validate_extended_benchmark.py}重新加载每个BPMN，校验文件摘要和消息迹数量，并重新计算异常参与方、MinAS和严格有效规则。154个新增输入均通过独立复核；在此基础上，又对6个合成规模的3类异常各抽取1个输入、对8个既有模型各抽取1个输入，共执行26次完整分析流水线冒烟测试，未发现结果不一致。

需要指出的是，新增154个输入均按照本文方法能够表达的异常类型设计，并在保留前经过严格有效性筛选。因此，这一数据组的154/154检出和修复结果用于证明实现对目标规则的覆盖能力、验证回归稳定性并开展规模实验，不能被解释为未知分布上的无偏准确率。本文仍使用原始67个案例输入的$65/67$可观测检出比例和$44/67$严格日志级修复比例回答实际案例中的有效性问题，并单独报告扩展基准结果，以避免由样本选择造成评价偏差。

\subsection{扩展规模性能实验}
\label{subsec:扩展规模性能实验}

规模性能实验使用\path{scripts/benchmark_extended_scale.py}完成。在正式计时前，先执行一次PM4Py模型加载作为预热；随后对6个合成BPMN的全部90个日志各执行一次完整分析。计时直接采用程序生成的\texttt{elapsed\_seconds}字段，包括BPMN规范化、流程树转换、消息模式分析、MinAS求解、修复候选生成、严格复验和输入摘要计算，不包括Python及PM4Py模块导入、JSON写入和终端输出。实验环境为Windows 10、AMD Ryzen 5 5600、Python 3.12.4和PM4Py 2.7.19.8。每个规模包含15次运行，分别覆盖5个消息缺失、5个顺序逆转和5个选择冲突输入。

\begin{table}[htbp]
\centering
\small
\renewcommand{\arraystretch}{1.25}
\caption{规模递增合成BPMN的预热分析性能}
\label{tab:规模递增合成BPMN性能}
\begin{tabularx}{\textwidth}{
    >{\centering\arraybackslash}X
    >{\centering\arraybackslash}X
    >{\centering\arraybackslash}X
    >{\centering\arraybackslash}X
    >{\centering\arraybackslash}X
    >{\centering\arraybackslash}X
    >{\centering\arraybackslash}X
}
\toprule
规模编号 & 参与方数 & 消息流数 & 日志数 & 检出/严格修复 & 平均耗时（秒） & P95（秒） \\
\midrule
Scale-1 & 5 & 12 & 15 & 15/15 & 0.0768 & 0.1061 \\
Scale-2 & 7 & 20 & 15 & 15/15 & 0.1252 & 0.1422 \\
Scale-3 & 9 & 32 & 15 & 15/15 & 0.2734 & 0.3560 \\
Scale-4 & 12 & 48 & 15 & 15/15 & 0.5806 & 0.6967 \\
Scale-5 & 16 & 72 & 15 & 15/15 & 2.3086 & 2.4601 \\
Scale-6 & 20 & 96 & 15 & 15/15 & 5.1961 & 5.4534 \\
\bottomrule
\end{tabularx}
\end{table}

如表\ref{tab:规模递增合成BPMN性能}所示，在参与方数量由5增加至20、消息流数量由12增加至96时，平均分析耗时由0.0768秒增加至5.1961秒。前4个规模的平均耗时均低于0.6秒；当消息流增加至72和96时，耗时分别上升至2.3086秒和5.1961秒。这一变化表明，当前实现在中小规模模型上开销较低，但包含更多并行块的较大模型会显著增加消息模式展示枚举、流程树转换和修复行为验证的计算量。最大模型的15次运行P95为5.4534秒，且15个输入均完成异常定位和严格修复复验，没有发生超时或枚举结果被静默截断。

从Scale-1到Scale-6，消息流数量增加8倍，而平均耗时约增加67.6倍，说明实验区间内的运行时间并非随消息流数量线性增长。该结果比仅使用原始工业案例更清楚地揭示了并行组合空间带来的性能代价，但仍不能直接作为算法渐近复杂度的证明。后续可以通过固定控制流结构、分别改变参与方数量、消息流数量和并行块数量的消融实验，进一步区分各因素对耗时和内存占用的影响。

